\subsection{2.2.1 Motivation relationaler Kernel im Ontophänomenalen Phasenraum}

\subsubsection{Einleitung}

Der Ontophänomenale Phasenraum (OPP) besitzt nach den bisherigen
Grundannahmen der PST zunächst keinerlei klassische Geometrie.

Insbesondere existieren im OPP zunächst keine

\begin{itemize}
\item kartesischen Koordinaten,
\item physikalischen Abstände,
\item Winkel,
\item Metriken,
\item Dimensionen
\end{itemize}

im üblichen Sinn.

Der OPP stellt daher keinen Raum im mathematischen Sinne einer
Mannigfaltigkeit mit vorgegebener Metrik dar.

Vielmehr bildet er ein rein ontologisches Relationssubstrat.

Die erste mathematische Aufgabe besteht deshalb nicht darin,
Koordinaten einzuführen, sondern die Beziehungen zwischen
ontophänomenalen Zuständen zu beschreiben.

Genau hierfür werden relationale Kernel eingeführt.

\bigskip

\subsubsection{Der relationale Grundgedanke}

Seien

\[
x_i,\;x_j \in \mathcal{O}
\]

zwei Zustände des OPP.

Zwischen diesen Zuständen existiert zunächst lediglich eine
Relation

\[
R_{ij}.
\]

Diese Relation besitzt zunächst keine Interpretation als Abstand.

Stattdessen beschreibt sie lediglich die Stärke der gegenseitigen
Verknüpfung.

Formal

\[
R :
\mathcal O \times \mathcal O
\longrightarrow
\mathbb R .
\]

Dabei gilt

\[
R_{ij}=R_{ji}.
\]

Die Relation ist somit zunächst symmetrisch.

Die gesamte Information des OPP wird in der Matrix

\[
R=
(R_{ij})
\]

gespeichert.

Diese Matrix stellt das eigentliche Fundament der PST dar.

Nicht die Punkte selbst sind fundamental,

sondern ihre Beziehungen.

Dies entspricht unmittelbar der relationalen Grundidee des
Ontophänomenalismus.

\bigskip

\subsubsection{Warum überhaupt Kernel?}

Eine beliebige Matrix wäre mathematisch möglich.

Sie besitzt jedoch keinerlei Struktur.

Die Physik zeigt dagegen überall eine sehr starke Lokalität.

Beispiele sind:

\begin{itemize}
\item Gravitationsfelder
\item elektromagnetische Felder
\item Quantenkorrelationen
\item statistische Korrelationen
\item Diffusionsprozesse
\end{itemize}

Alle besitzen gemeinsam:

\begin{quote}
Je ähnlicher zwei Zustände sind,
desto stärker beeinflussen sie sich.
\end{quote}

Ein Kernel beschreibt genau diese Idee.

Er ordnet jedem Paar

\[
(x_i,x_j)
\]

eine kontinuierliche Kopplungsstärke zu.

\bigskip

\subsubsection{Minimalforderungen an einen PST-Kernel}

Aus physikalischer Sicht ergeben sich mehrere
Minimalforderungen.

\paragraph{1. Positivität}

Die Kopplungsstärke soll nicht negativ werden.

\[
R_{ij}\ge0.
\]

Negative Kopplungen können später eingeführt werden,
werden jedoch zunächst ausgeschlossen.

\paragraph{2. Symmetrie}

Da zunächst keine ausgezeichnete Richtung existiert,

muss gelten

\[
R_{ij}=R_{ji}.
\]

\paragraph{3. Maximale Selbstkopplung}

Ein Zustand besitzt stets maximale Ähnlichkeit mit sich selbst.

Daher

\[
R_{ii}=1.
\]

Dies normiert automatisch den gesamten Kernel.

\paragraph{4. Monotonie}

Mit wachsender Trennung soll die Kopplung
abnehmen.

Ist

\[
d_1<d_2,
\]

so soll gelten

\[
R(d_1)>R(d_2).
\]

Dies beschreibt mathematisch die intuitive Forderung
lokaler Wechselwirkungen.

\paragraph{5. Glattheit}

Kleine Änderungen der Relation sollen kleine Änderungen
der Kopplung erzeugen.

Der Kernel soll daher stetig sein.

\bigskip

\subsubsection{Der Kernel als Projektion der Ontologie}

Im OP existieren zunächst keine physikalischen Abstände.

Die Funktion

\[
R(d)
\]

ist daher nicht als Abstandsgesetz zu verstehen.

Sie beschreibt vielmehr,

wie stark zwei ontophänomenale Zustände
gegenseitig miteinander gekoppelt sind.

Erst spätere Projektionen erzeugen aus dieser
relationalen Information eine effektive Geometrie.

Die Reihenfolge lautet daher

\[
\boxed{
\text{Relation}
\;\longrightarrow\;
\text{Kernel}
\;\longrightarrow\;
\text{Spektrum}
\;\longrightarrow\;
\text{Geometrie}
}
\]

nicht umgekehrt.

Dies stellt einen fundamentalen Unterschied
zur klassischen Differentialgeometrie dar.

\bigskip

\subsubsection{Vom Kernel zur Gram-Matrix}

Aus dem Kernel wird anschließend eine
Gram-Matrix konstruiert.

Sie besitzt allgemein die Form

\[
K=f(R).
\]

In den bisherigen PST-Simulationen wurde unter anderem

\[
K_{ij}
=
R_{i0}
+
R_{j0}
-
R_{ij}
\]

verwendet.

Diese Konstruktion besitzt mehrere Vorteile:

\begin{itemize}
\item sie ist symmetrisch,
\item sie erlaubt Spektralanalysen,
\item sie erzeugt Eigenmoden,
\item sie ermöglicht Dimensionsschätzungen,
\item sie ist numerisch stabil.
\end{itemize}

Die Eigenwerte von

\[
K
\]

enthalten somit die gesamte geometrische Information
der jeweiligen Projektion.

\bigskip

\subsubsection{Warum verschiedene Kernel getestet werden müssen}

Eine zentrale offene Frage der PST lautet:

\begin{quote}
Ist die beobachtete Raumzeit eine Eigenschaft des OPP selbst
oder lediglich eine Eigenschaft einer speziellen Projektion?
\end{quote}

Um diese Frage beantworten zu können,
muss untersucht werden,

wie empfindlich die entstehende Geometrie
gegenüber der Wahl des Kernels ist.

Ist beispielsweise ausschließlich ein bestimmter
Kernel in der Lage,

eine stabile effektive Vierdimensionalität
zu erzeugen,

dann besitzt dieser Kernel möglicherweise
eine fundamentale physikalische Bedeutung.

Entstehen dagegen ähnliche Ergebnisse
für viele unterschiedliche Kernel,

spricht dies für eine universelle Eigenschaft
des OPP.

Die Kernelwahl wird dadurch selbst
zu einer überprüfbaren physikalischen Hypothese.

\bigskip

\subsubsection{Erkenntnistheoretische Bedeutung}

Die bisherigen numerischen Untersuchungen legen nahe,

dass die beobachtete Raumzeit nicht unmittelbar
im OPP vorhanden ist.

Vielmehr scheint sie als Eigenschaft
bestimmter Projektionen aufzutreten.

Der Kernel übernimmt dabei die Rolle eines
Projektionsoperators.

Er bestimmt,

welche Aspekte der hochdimensionalen relationalen Struktur
überhaupt sichtbar werden.

Damit erhält der Kernel eine doppelte Bedeutung:

\begin{enumerate}
\item mathematisch beschreibt er Relationen,
\item erkenntnistheoretisch beschreibt er die Art,
wie Beobachter den OPP wahrnehmen.
\end{enumerate}

Dies passt unmittelbar zur Grundidee des
Ontophänomenalismus:

Nicht die Welt besitzt notwendigerweise
eine vierdimensionale Struktur,

sondern unsere Beobachtung könnte genau jene
Projektionsvorschrift realisieren,
welche aus dem hochdimensionalen OPP
eine stabile emergente Raumzeit entstehen lässt.

